\section{Results}
\label{sec:results}

%
% 3.1
%

\subsection{Baseline Attack Execution (Control Group)}
The baseline testing phase confirmed the presence of the Scanner Gap vulnerability within both unmitigated container environments (Docker and Apptainer). In the control group trials, the fileless loader executed the payload in memory across both runtimes, successfully evading all static signature checks. Upon execution, the payload established an active reverse shell TCP connection back to the command and control (C2) listener, as documented in the live session log in Listing \ref{lst:c2_connection} for the Docker trial; the Apptainer trial produced functionally identical results, differing only in the target C2 gateway address (\texttt{127.0.0.1} due to the shared host network namespace) and hostname mapping. These connections granted the remote attacker full interactive root shell access inside the running container namespaces without writing any executable binary to the persistent storage layers.


\begin{lstlisting}[language=bash, caption={Terminal output demonstrating successful payload execution and reverse shell connection in Docker.}, label={lst:c2_connection}]
[PAYLOAD] Exfil proof file: /tmp/containment_leak_PROOF.txt
[PAYLOAD] --- Initiating reverse shell to 172.17.0.1:4444 ---
[PAYLOAD] FILELESS PAYLOAD EXECUTION CONFIRMED
[PAYLOAD] Timestamp : 2026-06-22T06:38:59 UTC
[PAYLOAD] PID       : 86
[PAYLOAD] PPID      : 80
[PAYLOAD] TARGET HOSTNAME : 2f3d36941775
[PAYLOAD] Kernel : Linux 6.12.74+deb13+1-amd64 x86_64
\end{lstlisting}

The evasion mechanics rely entirely on anonymous memory allocation. As shown in Listing \ref{lst:baseline_maps}, the active process's memory layout (accessed via \texttt{/proc/self/maps}) displays the payload's ELF segments mapped to the anonymous file descriptor \texttt{/memfd: (deleted)} on device \texttt{00:01} with inode \texttt{1079}. This layout contains multiple distinct entry ranges for the same memory file. These five entries represent individual program segments mapped by the kernel's Memory Management Unit to separate Virtual Memory Areas (VMAs). Although they share three categories of permissions, \texttt{r--p} for read-only headers and constants, \texttt{r-xp} for executable machine instructions, and \texttt{rw-p} for writeable variables, they are partitioned into five distinct memory ranges to align with 4KB page boundaries and support modern binary security hardening such as relocation read-only (RELRO). 

\begin{lstlisting}[language=bash, caption={Memory mapping of the fileless payload demonstrating the Scanner Gap.}, label={lst:baseline_maps}]
[PAYLOAD] --- /proc/self/maps (first 15 lines) ---
[PAYLOAD] [FILELESS] >> 00400000-00401000 r--p 00000000 00:01 1079   /memfd: (deleted)
[PAYLOAD] [FILELESS] >> 00401000-004a2000 r-xp 00001000 00:01 1079   /memfd: (deleted)
[PAYLOAD] [FILELESS] >> 004a2000-004ce000 r--p 000a2000 00:01 1079   /memfd: (deleted)
[PAYLOAD] [FILELESS] >> 004ce000-004d3000 r--p 000cd000 00:01 1079   /memfd: (deleted)
[PAYLOAD] [FILELESS] >> 004d3000-004d5000 rw-p 000d2000 00:01 1079   /memfd: (deleted)
[PAYLOAD] --- end of maps ---
[PAYLOAD] HYPOTHESIS CONFIRMED: Payload executed with NO on-disk ELF.
\end{lstlisting}

\clearpage 

The major and minor device numbers (\texttt{00:01}) indicate that the backing file resides on an anonymous virtual memory mount (\texttt{tmpfs}/\texttt{shmem}) rather than a physical block device, while the inode resides entirely within the kernel's RAM-backed allocation tables. Because the memory space is unlinked from the directory tree, as indicated by the \texttt{(deleted)} suffix, it is inaccessible to standard file path traversals. Consequently, since traditional endpoint protection agents rely on intercepting physical disk write events or scanning persistent block storage, executing the payload directly from this anonymous memory-resident structure generates zero filesystem writes. The absence of physical artifacts prevents static signature detection. This lack of an on-disk file path confirms the hypothesis that the payload executed without a physical ELF file; because the process's symlink target (\texttt{/proc/self/exe}) references an unlinked virtual memory descriptor rather than a persistent directory node, the entire execution cycle is isolated within RAM. Table \ref{tab:scanner_gap_forensics} details the audit of the environment during this baseline test.


\begin{table}[htbp]
\centering
\caption{Forensic Audit of the Static Scanner Gap.}
\label{tab:scanner_gap_forensics}
\begin{tabular}{@{}p{0.3\linewidth}p{0.25\linewidth}p{0.35\linewidth}@{}}
\toprule
\textbf{Metric} & \textbf{Status} & \textbf{Evidence} \\ \midrule
File System Artifacts & Null & No inode mapping found \\
Volatility Inspection & Active & Bytes resident in RAM \\
Signature Detection & Negative & No pattern match for ELF \\ \bottomrule
\end{tabular}
\end{table}

%
% 3.2
%

\subsection{eBPF Mitigation Performance (Docker)}
The Docker-based eBPF mitigation module intercepted the offensive execution chain. When the fileless loader invoked \texttt{memfd\_create} and \texttt{execve}, the kernel tracepoints triggered the stateful BPF filter. The stateful eBPF monitor correlated the system call chain. The kernel ring buffer recorded the detection event within the five-second correlation window, listing the execution cost alongside the PID metadata (Listing \ref{lst:ebpf_telemetry}). This correlation triggered the immediate termination of the execution sequence. The kernel returned an \texttt{-EPERM} error code to the user-space process (Listing \ref{lst:ebpf_denial}).

The monitor blocked 100\% of execution attempts. The enforcement hook ran in 832 nanoseconds; this wall-time cost was measured directly within the kernel space by bracketing the BPF execution block with the \texttt{bpf\_ktime\_get\_ns()} helper function. 


\begin{lstlisting}[language=bash, caption={eBPF telemetry recording the correlated system call violation and LSM hook execution latency.}, label={lst:ebpf_telemetry}]
[LSM] memfd->exec block active (window=5s, comm=loader)
[BLOCK] pid=31738 uid=0 comm=loader delta_ns=703570 lsm_exec_ns=832
\end{lstlisting}

\begin{lstlisting}[language=bash, caption={Loader process termination resulting from the LSM enforcement.}, label={lst:ebpf_denial}]
[LOADER] Step 2: memfd_create ... OK (fd=3, /proc/self/fd/3 -> /memfd: )
[LOADER] Step 3b: memfd sealed (F_SEAL_WRITE|SHRINK|GROW|SEAL) ... OK
[LOADER] Step 4: fexecve     ... calling (no return on success)
[LOADER] fexecve: Operation not permitted
\end{lstlisting}

%
% 3.3
%

\subsection{Seccomp Mitigation Performance (Apptainer)}
The Seccomp structural restriction neutralized the threat during container instantiation. The kernel evaluated the system calls against the declarative JSON configuration policy synchronously upon entry. The Apptainer configuration blocked 100\% of execution attempts. Because the kernel enforces the policy during the system call entry path, rejection latency averaged 292.82 nanoseconds. Listing \ref{lst:seccomp_denial} contains the loader error output verifying the kernel-level blocking of the initialization sequence.

\begin{lstlisting}[language=bash, caption={Loader process termination resulting from the Seccomp enforcement.}, label={lst:seccomp_denial}]
[LOADER] ============================================================
[LOADER] Fileless Loader
[LOADER] Technique : memfd_create + fexecve (T1027.002)
[LOADER] Payload   : 961928 bytes (XOR-encoded, key=0xAB)
[LOADER] ============================================================
[LOADER] Step 1: XOR decode  ... OK (961928 bytes)
[LOADER] memfd_create: Operation not permitted
\end{lstlisting}

% 
% 3.4
% 

\subsection{Comparative Performance Matrix}
Table \ref{tab:comparative_summary} synthesizes the experimental results from the attack and both defensive paradigms. 

\begin{table}[htbp]
\centering
\caption{Comparative analysis of baseline execution versus applied mitigation strategies.}
\label{tab:comparative_summary}
\begin{tabular}{@{}p{0.15\linewidth}p{0.22\linewidth}p{0.18\linewidth}p{0.15\linewidth}p{0.22\linewidth}@{}}
\toprule
\textbf{Environment} & \textbf{Loader Result} & \textbf{C2 Connection} & \textbf{Enforcement Latency} & \textbf{Key Forensic Evidence} \\ \midrule
Baseline (Control) & Unrestricted Execution & Established & 0 ns & \texttt{/memfd: (deleted)} mapped in memory \\
eBPF (Docker) & \texttt{fexecve: -EPERM} & Denied & 832 ns & eBPF ring buffer PID correlation \\
Seccomp (Apptainer) & \texttt{memfd\_create: -EPERM} & Denied & 292.82 ns & \texttt{memfd\_create: Operation not permitted} \\ \bottomrule
\end{tabular}
\end{table}